\documentclass{ctexbook}
\usepackage{EC}
\usepackage[
    backend=biber,
    style=chem-acs,
]{biblatex}
\addbibresource{ref.bib}
\begin{document}
\section{氮及其化合物}
\subsection{氮的卤化物与卤氧化物}
1811年,P.L.Dulong历经千辛万苦制得了\ce{NCl3}(他为此付出了三根手指和一只眼睛的代价).令人费解的是\ce{N}最稳定的卤化物\ce{NF3}直到115年后才被制得.
\subsubsection{氮的氟化物}
\paragraph{\ce{NF3}}
纯的\ce{NF3}首先是利用电解\ce{NH4F/HF}熔融物而制备的;另一个可用的方法是在\ce{Cu}催化剂上对\ce{NH3}进行控制氟化:
\begin{center}
    \ce{NH4F + 2HF ->T[熔融][电解] NF3 + 3H2}\\
    \ce{4NH3 + 3F2 ->T[\ce{Cu}] NF3 + 3NH4F}
\end{center}
\begin{substance}[\ce{NF3}]
    三氟化氮,化学式为\ce{NF3},是无色无臭的稳定气体,熔点为$-206.8\tc$,沸点为$-129.0\tc$.
\end{substance}
\indent\ce{NF3}分子为三角锥形,键角为$102.5^\circ$,偶极矩异常地小(仅为\ce{NH3}的六分之一).这是由于\ce{N-F}键矩方向与孤对电子方向相反所致.\\
\indent 在常温下,\ce{NF3}是不活泼的.升高温度时\ce{NF3}可以作为氟化剂:
\begin{center}
    \ce{2NF3 + 2Cu -> N2F4 + CuF2}
\end{center}
尽管\ce{NF3}是稳定的,但它仍然能形成一些四面体构型的离子:
\begin{center}
    \ce{NF3 + F2 + SbF5 ->T[$200\tc$][100 atm][NF4]+[SbF6]-}\\
    \ce{2NF3 + O2 ->T[放电][$-196\tc$] 2ONF3}\\
    \ce{3NOF + 2IrF6 ->T[$20\tc$] ONF3 + 2[NO]+[IrF6]-}
\end{center}
\ce{NOF3}中有着明显短的\ce{N-O}键(115.8 pm)和明显长的\ce{N-F}键(143.1 pm).
\paragraph{\ce{N2F4}}
\ce{N2F4}除了用\ce{Cu}还原\ce{NF3}得到外还可以通过以下方式得到:
\begin{center}
    \ce{CO(NH2)2 + 2F2 ->T[\ce{N2/F2}] H2NCONF2 + 2HF}\\
    \ce{NH2CONF2 ->T[浓\ce{H2SO4}] NF2H + HOCN}\\
    \ce{2NF2H + NaOCl -> N2F4 + H2O + NaCl}
\end{center}
\begin{substance}[\ce{N2F4}]
    四氟化二氮,化学式为\ce{N2F4},是无色气体,熔点为$-164.5\tc$,沸点为$-73\tc$.
\end{substance}
\indent\ce{N2F4}是较为活泼的气体,可以作为氟化剂与很多物质发生反应,例如:
\begin{center}
    \ce{10Li + N2F4 -> 4LiF + 2Li3N}
\end{center}
此外,\ce{N2F4}还能离解产生\ce{F2N.}自由基.在低压下快速冷冻\ce{N2F4}就能得到暗蓝色的固体.事实上,\ce{N2F4}的很多化学性质也与\ce{F2N.}自由基有关,例如:
\begin{center}
    \ce{N2F4 + 2PhSH -> PhSSPh + 2NF2H}\\
    \ce{N2F4 + 2NO -> 2F2NNO}
\end{center}
\paragraph{\ce{N2F2}}最好的制备\ce{N2F2}的方法如下:
\begin{center}
   \ce{ KF + NF2H ->T[$-80\tc$][then $20\tc$] KHF2 + N2F2}
\end{center}
制得的\ce{N2F2}一般是顺式和反式异构体的混合物.
\begin{substance}[\ce{N2F2}]
    二氟化二氮,化学式为\ce{N2F2},是无色气体.\\
    cis-\ce{N2F2}的熔点低于$-195\tc$,沸点为$-105.7\tc$;trans-\ce{N2F2}的熔点为$-172\tc$,沸点为$-111.4\tc$.
\end{substance}
下面是两种异构体的键长数据.
\begin{table}[H]
    \centering
    \begin{tabular}{|c|c|c|}
        \hline
        键长/pm&\ce{N=N}&\ce{N-F}\\\hline
        顺式\ce{N2F2}&120.9&140.9\\\hline
        反式\ce{N2F2}&122.4&139.8\\\hline
    \end{tabular}
\end{table}
造成键长差异的原因主要是\ce{N}的孤对电子对$\sigma^\ast\left(\ce{N-F}\right)$轨道的超共轭作用,这一作用削弱\ce{N-F}键而加强\ce{N=N}键,也为分子提供了额外的稳定化作用.因此,顺式\ce{N2F2}在热力学上更加稳定,在气相的平衡中也占据主要.\\
\indent 然而,顺式\ce{N2F2}在动力学上似乎更加活泼.顺式\ce{N2F2}在玻璃容器内存放两周即完全反应:
\begin{center}
    \ce{SiO2 + 2N2F2 -> SiF4 + 2N2O}
\end{center}
反式\ce{N2F2}则可以稳定存放在玻璃容器中.此外,只有顺式\ce{N2F2}能脱去\ce{F-}形成\ce{N2F+}离子.这也是\ce{AsF5}分离两者的原理.
\subsubsection{氮的其余卤化物}
\paragraph{\ce{NCl3}}
除去氟化物外,\ce{NCl3}几乎是氮唯一稳定的卤化物.
\begin{substance}[\ce{NCl3}]
    三氯化氮,化学式为\ce{NCl3},是黄色油状液体,熔点$-40\tc$,沸点$70\tc$.
\end{substance}
\indent\ce{NCl3}可以通过电解酸性\ce{NH4Cl}溶液制得.它主要用于漂白与杀菌.\ce{NCl3}容易水解:
\begin{center}
    \ce{NCl3 + 3H2O -> NH3 + 3HOCl}\\
\end{center}
在碱性溶液中可以与\ce{NaClO2}反应制取\ce{ClO2}:
\begin{center}
    \ce{NCl3 + 3H2O + 6NaClO2 -> 6ClO2 + 3NaCl + 3NaOH + NH3}\\
    \ce{2NCl3 + 6NaClO2 -> 6ClO2 + 6NaCl + N2}
\end{center}
\paragraph{\ce{NBr3}和\ce{NI3}}
这两种化合物都十分的不稳定.
\subsubsection{氮的卤氧化物}
\paragraph{亚硝酰卤化物\ce{NOX}}
亚硝酰卤化物都是活泼气体.它们的物理性质和结构可列举如下:
\begin{table}[H]\centering
    \begin{tabular}{cccc}
        \hline
        性质    &\ce{NOF}   &\ce{NOCl}  &\ce{NOBr} \\\hline
        颜色    &无色       &橙黄色     &红色 \\
        熔点/$\tc$  &$-132.5$   &$-59.6$    &$-56$ \\
        沸点/$\tc$  &$-59.9$   &$6.4$    &$\sim0$ \\
        键角    &$110^\circ$    &$113^\circ$    &$117^\circ$\\
        \ce{N=O}键长/pm &$113$  &$114$  &$115$\\\hline
    \end{tabular}
\end{table}
它们可以通过以下方法制备:
\begin{center}
    \ce{2NO + X2 -> 2NOX}\ \ \ \ce{X=F,Cl,Br}\\
    \ce{NO + AgF2 -> NOF + AgF}\\
    \ce{N2O4 + KCl -> NOCl + KNO3}
\end{center}
\ce{NOX}在\ce{X-}受体的存在下可以裂解产生\ce{NO+},例如
\begin{center}
    \ce{NOF + BF3 -> [NO]+[BF4]-}\\
    \ce{NOCl + AlCl3 -> [NO]+[AlCl4]-}
\end{center}
\paragraph{硝酰卤化物\ce{NO2X}}
硝酰氟化物和氯化物\ce{XNO2},如同它的亚硝酰类似物,也是活泼气体;它们皆为平面形,和等电子的硝酸根\ce{NO3-}类似.它们的物理性质列举如下.
\begin{table}[H]\centering
    \begin{tabular}{ccc}
        \hline
        性质    &\ce{NO2F}   &\ce{NO2Cl} \\\hline
        熔点/$\tc$  &$-166$   &$-145$ \\
        沸点/$\tc$  &$-72.5$   &$-15.9$ \\
        \ce{X-N-O}键角    &$118^\circ$    &$115^\circ$ \\
        \ce{N-O}键长/pm &$123$  &$120$ \\\hline
    \end{tabular}
\end{table}
\ce{XNO2}的反应常常与\ce{XNO}类似.例如形成\ce{NO2+}的盐:
\begin{center}
    \ce{NO2F + BF3 -> [NO2]+[BF4]-}\\
    \ce{NO2Cl + SbCl5 -> [NO2]+[SbCl6]-}
\end{center}
\ce{NO2Cl}在水和液氨中的溶剂解也生成不同的产物.
\begin{center}
    \ce{ClNO2 + H2O -> HNO3 + HCl}\\
    \ce{ClNO2 + 2NH3 -> ClNH2 + NH4NO2}
\end{center}
造成这一差异的原因是在水中\ce{HNO2}会被\ce{HClO}氧化,而在液氨中则不发生类似的反应.
\subsection{氮的氧化物}
在众多的元素中间,氮是唯一的能形成多达(至少)7种氧化物分子的元素.所有的氮的氧化物在热力学上都是不稳定的,可分解为\ce{N2}和\ce{O2}.
\subsubsection{\ce{N2O}}
\begin{substance}[\ce{N2O}]
    一氧化二氮,俗称笑气,化学式为\ce{N2O},是无色有甜味的气体.\ce{N2O}的熔点为$-90.86\tc$,沸点为$-88.48\tc$,可以少量地溶于水中.
\end{substance}
\paragraph{\ce{N2O}的结构}
作为\ce{CO2}的等电子体,\ce{N2O}是直线型分子.理论计算表明其中\ce{N=N}键级为$2.73$,\ce{N-O}键级为$1.61$.
\paragraph{\ce{N2O}的制备}
小心地加热熔融的\ce{NH4NO3}可以使其热分解得到\ce{N2O}:
\begin{center}
    \ce{NH4NO3 ->T[$\Delta$] N2O + 2H2O}
\end{center}
这个反应的历程比较复杂.此外还有可以采用的方法为
\begin{center}
    \ce{HNO2 + NH2OH ->T[aq] N2O + 2H2O}\\
    \ce{HNO3 + HN3 ->T[aq] N2O + N2 + H2O}\\
    \ce{H2NNO2 ->T[$\Delta$] N2O + H2O}\ \ \ \ \ \ce{H2N2O2 ->T[$\Delta$] N2O + H2O}
\end{center}
尽管加热连二次硝酸\ce{H2N2O2}确实可以使其脱水得到\ce{N2O},但我们仍然不认为\ce{N2O}是连二次硝酸酐,因为\ce{N2O}溶于水并不能再次形成\ce{H2N2O2}.
\paragraph{\ce{N2O}的化学性质与反应}
室温下\ce{N2O}较不活泼,例如它不能和卤素,碱金属甚至臭
氧反应.在较高温度时它可以作为氧化剂氧化一些物质,同时生成\ce{N2}.\\
\indent \ce{N2O}的一个重要应用是制备叠氮化物.我们已经在前面介绍过.\\
\indent \ce{N2O}最大的工业应用是作为火箭推进剂和搅打起泡冰激淋(或奶油)的充气剂,这是依靠在压力下\ce{N2O}在植物油中的溶解度,以及它在低浓度时无毒无味.\ce{N2O}也大量用作麻醉剂.\\
\indent \ce{N2O}得名为笑气是因为吸入它使人欢欣并且发笑.然而,过量吸入\ce{N2O}会导致神志不清,对神经系统造成危害\footnote{\ce{N2O}能不可逆地氧化钴胺素的钴核并使甲基钴胺素失活,间接诱发维生素B12缺乏,造成神经损伤.},同时有一定的成瘾性.因此,\ce{N2O}在我国被列入《危险化学品目录》,其销售和购买需要严格监管.\\
\indent 此外,\ce{N2O}也是一种温室气体.
\subsubsection{\ce{NO}与\ce{N2O2}}
\begin{substance}[\ce{NO}]
    一氧化氮,化学式为\ce{NO},是无色无味气体,熔点为$-164\tc$,沸点为$-152\tc$.
\end{substance}
\paragraph{\ce{NO}的结构}
作为最简单的具有单电子的分子,人们对\ce{NO}的电子结构进行了详细的研究,这也推动了分子轨道理论的发展.
\paragraph{\ce{NO}的制备}
在工业上,\ce{NO}是由\ce{NH3}的催化氧化得到的,这在前面已经有所提及.\\
\indent 实验室中制取\ce{NO}可以通过以下几种方式:
\begin{center}
    \ce{2KNO2 + 2KI + 2H2SO4 -> 2NO + 2K2SO4 + H2O + I2}\\
    \ce{KNO2 + K4[Fe(CN)6] + 2CH3COOH -> NO + K3[Fe(CN)6] + H2O + 2CH3COOK}\\
    \ce{6NaNO2 + 3H2SO4 -> 4NO + 2HNO3 + 2H2O + 3Na2SO4}\\
    \ce{3KNO2 + KNO3 + Cr2O3 -> 4NO + 2K2CrO4}
\end{center}
\paragraph{\ce{NO}的二聚形成\ce{N2O2}}
正常情况下,\ce{NO}可以少量地二聚形成\ce{O=N-N=O}分子(先前也有人认为二聚体是矩形环状结构的,不过这被晶体的衍射数据所否定).\\
\indent 在极性物质或Lewis酸存在时冷凝\ce{NO},则可以得到红色的不对称二聚体\ce{O=N-O=N}.
\paragraph{\ce{NO}的化学性质}
\ce{NO}具有较为活泼的化学性质.\\
\indent 首先,\ce{NO}易被\ce{O2}氧化为\ce{NO2}:
\begin{center}
    \ce{2NO + O2 -> 2NO2}
\end{center}
在空气中的\ce{NO}会很快地被氧化为\ce{NO2},从而显现红棕色.这一反应还有一个有趣之处在于它是三级反应,并且表观活化能为负值.这意味着温度的升高会使得反应速率减慢.事实上,这一反应的机理如下:
\begin{center}
    \ce{2NO <=>T[$K$][快速平衡] N2O2}\\
    \ce{N2O2 + O2 ->T[$k$][慢反应] 2NO2}
\end{center}
因此反应的速率方程为$v=Kk\con{NO}^2\con{O2}$.当温度升高时,$K$减小的程度比$k$增加的程度更大,因此反应的总体速率呈现下降趋势.\\
\indent 前面已经提到过,\ce{NO}容易与卤素反应生成对应的亚硝酰卤化物.\\
\indent 鉴于\ce{NO}并没有对应价态的稳定阴离子,因此\ce{NO}在碱性条件下容易发生歧化反应:
\begin{center}
    \ce{2Na2O + 4NO ->T[$100\tc$] 2NaNO2 + Na2N2O2}\\
    \ce{2MOH + 4NO -> 2MNO2 + N2O + H2O}\\
    \ce{4MOH + 6NO -> 4MNO2 + N2 + 2H2O}
\end{center}
\paragraph{\ce{NO}的配合物}
\ce{NO}易和许多过渡金属化合物反应形成亚硝酰配合物,而且在其他的含有氧-氮物种的反应中也经常形成这类配合物.经典的例子是检验\ce{NO3-}或\ce{NO2-}棕色环实验:
\begin{center}
    \ce{3[Fe(H2O)6]^2+ + NO3- + 4H+ -> 2[Fe(H2O)6]^3+ + [Fe(H2O)5(NO)]^2+ + 3H2O}\\
    \ce{2[Fe(H2O)6]^2+ + NO2- + 2H+ -> [Fe(H2O)6]^3+ + [Fe(H2O)5(NO)]^2+ + 2H2O}
\end{center}
取待测溶液于试管中,加入适量\ce{FeSO4}溶液.缓慢沿管壁加入浓硫酸.如果在浓硫酸与溶液的界面观察到了棕色环,说明溶液中含有\ce{NO3-}或\ce{NO2-}.使用醋酸代替浓硫酸,则\ce{NO3-}不发生此反应,而\ce{NO2-}则正常.\\
\indent 通常而言,\ce{NO}在配合物中有多样的几何构型.直线型通常意味着\ce{NO}以\ce{NO+}的形式配位,作为三电子配体\footnote{形式上可以这么认为,但实际上电荷分布并非如此,\ce{NO}是通过轨道作用提供三个电子的,并非预先发生了氧化还原反应,否则会导致中心金属元素的异常低价态.}.折线形通常意味着\ce{NO}以\ce{NO-}的形式配位,作为单电子供体.\\
\indent 由于上述原因,并且\ce{NO+}与\ce{CO}为等电子体,因此简单羰基化合物中的三个\ce{CO}可以替换为两个\ce{NO}.例如
\begin{center}
    \ce{Fe(CO)5 + 2NO -> Fe(NO)2(CO)2 + 3CO}\\
    \ce{Cr(CO)6 + 4NO -> Cr(NO)4 + 6CO}
\end{center}
其它\ce{NO}的配合物将在以后提及.
\subsubsection{\ce{N2O3}}
\begin{substance}[\ce{N2O3}]
    三氧化二氮,化学式为\ce{N2O3},是一种浅蓝色固体,熔点为$-100.1\tc$,熔化后为深蓝色液体.\ce{N2O3}对热不稳定,容易分解为\ce{NO}和\ce{NO2}.
\end{substance}
\paragraph{\ce{N2O3}的结构}
\ce{N2O3}中有着异常长的\ce{N-N}键.这暗示了它离解为两种自由基\ce{NO}和\ce{NO2}的明显趋势.
\paragraph{\ce{N2O3}的化学性质}
前面已经说到\ce{N2O3}很容易分解.随着温度的升高,混合物将变为绿色(由\ce{NO2}的红棕色和\ce{N2O3}的蓝色混合而成).早先的研究总是认为\ce{NO}的液体或固体是淡蓝色的,但事实上这也是由于其中混有的少量\ce{N2O3}所显现的强烈的蓝色所致.\\
\indent 制取\ce{N2O3}的最佳方法是冷凝\ce{NO}和\ce{NO2}的等比例混合物.向\ce{NO}中混入预定比例的\ce{O2}以就地合成\ce{NO2}也是一个可采取的方法.\\
\indent 作为亚硝酸盐的酸酐,\ce{N2O3}和水反应生成\ce{HNO2},而与碱反应几乎定量地生成\ce{NO2-}.此外,在足够浓的酸中,它也可以被转化为\ce{NO+}:
\begin{center}
    \ce{N2O3 + 3H2SO4(conc) -> 2[NO]+[HSO4]- + [H3O]+[HSO4]-}
\end{center}
\subsubsection{\ce{NO2}与\ce{N2O4}}
虽然大多数时候,分别测定平衡混合物中各物质的状态是可行的,但由于\ce{NO2}与\ce{N2O4}处于极为快速的平衡中,因此大多数时候我们描述的都是平衡混合物的性质.
\begin{substance}[\ce{NO2}]
    二氧化氮,化学式为\ce{NO2},为红棕色气体,凝固点为$-11.2\tc$,沸点为$21.5\tc$.固态和液态的\ce{NO2}中实际上几乎全部为无色的\ce{N2O4}.
\end{substance}
\paragraph{\ce{NO2}和\ce{N2O4}的结构}
相比\ce{NO},\ce{NO2}中的单电子更偏向于定域在\ce{N}原子上(即SOMO轨道为\ce{N}的$\text{sp}^2$轨道).因此,\ce{NO2}发生二聚作用相对更容易一些.\\
\indent 二聚体\ce{N2O4}为平面型结构,其中的\ce{N-N}键也格外的长.
\paragraph{\ce{NO2}和\ce{N2O4}的制备}
制备\ce{NO2}最好的办法是将\ce{Pb(NO3)2}热分解的产物分级蒸馏后得到:
\begin{center}
    \ce{2Pb(NO3)2 ->T[$400\tc$] 2PbO + 4NO2 + O2}
\end{center}
其它的办法(相较而言不佳)有:
\begin{center}
    \ce{2HNO3 + SO2 -> N2O4 + H2SO4}\\
    \ce{4HNO3 + P4O10 -> 2N2O4 + O2 + 4HPO3}\\
    \ce{NOCl + AgNO3 -> N2O4 + AgCl}
\end{center}
可以看出上面的反应都可以看做先得到\ce{N2O5},再分解生成\ce{N2O4}.
\paragraph{\ce{N2O4}作为溶剂}
\ce{N2O4}作为一种优良的非水溶剂得到了大量而广泛的研究.它的独特应用是制备各种金属的无水硝酸盐.\\
\indent 在极性物质存在的形况下,\ce{N2O4}倾向于按照以下的方式解离:
\begin{center}
    \ce{N2O4 <=> NO+ + NO3-}
\end{center}
这一溶剂体系内也可以发生类似于酸碱反应的反应:
\begin{center}
    \ce{NOCl + AgNO3 -> AgCl + N2O4}\\
    \ce{2NOCl + Sn -> SnCl2 + 2NO}\\
    \ce{2EtNH3NO3 + 2N2O4 + Zn -> [EtNH3]2[Zn(NO3)4] + 2NO}\\
    \ce{CaO + 2N2O4 -> Ca(NO3)2 + N2O3}
\end{center}
这类反应为制取无水硝酸盐提供了一个很好的办法.特别是当使用金属的碘化物时,由于\ce{NOI}迅速分解,因此减少了副反应的发生,例如
\begin{center}
    \ce{TiI4 + 4N2O4 -> Ti(NO3)4 + 4NO + 2I2}
\end{center}
\indent 与\ce{N2O4}倾向于裂解为\ce{NO+}和\ce{NO3-}相比,裂解为\ce{NO2+}和\ce{NO2-}的情形则要少得多.不过,\ce{N2O4}与\ce{BF3}反应可以生成下面两种加合物:
\begin{center}
    \ce{N2O4 + BF3 -> [NO2]+[ONOBF3]-}\\
    \ce{N2O4 + 2BF3 -> [NO2]+[F3BONOBF3]-}
\end{center}
\subsubsection{\ce{N2O5}}
\begin{substance}[\ce{N2O5}]
    五氧化二氮,化学式为\ce{N2O5},是无色柱状晶体,易溶于水并生成硝酸.固体\ce{N2O5}于$32.4\tc$升华.
\end{substance}
\paragraph{\ce{N2O5}的结构}
气相中的\ce{N2O5}分子呈现平面型结构.\\
\indent 固相的\ce{N2O5}实际上是\ce{[NO2]+[NO3]-},其中可以分别找到分立的两种离子.
\begin{figure}[H]
    \centering\includegraphics[scale=0.125]{figure/N2O5s.eps}\caption{固体\ce{N2O5}的晶体结构}
\end{figure}
\paragraph{\ce{N2O5}的制备与反应}
\ce{N2O5}是硝酸的酸酐,可以通过低温下小心地用\ce{P4O10}对\ce{HNO3}脱水得到:
\begin{center}
    \ce{4HNO3 + P4O10 -> 2N2O5 + 4HPO3}
\end{center}
\indent \ce{N2O5}是一种强氧化剂,可以与很多物质发生反应:
\begin{center}
    \ce{N2O5 + Na -> NaNO3 + NO2}\\
    \ce{N2O5 + I2 -> I2O5 + N2}
\end{center}
\indent 与\ce{N2O4}类似,\ce{N2O5}在极性溶剂(例如各种无水强酸)中易解离成离子.这为制取\ce{NO2+}的盐提供了一个好办法:
\begin{center}
    \ce{N2O5 + 3H2SO4 -> 2[NO2]+[HSO4]- + [H3O]+[HSO]4-}\\
    \ce{N2O5 + HSO3F -> [NO2]+[SO3F]- + HNO3}
\end{center}
\subsection{氮的含氧酸及其盐}
\subsubsection{连二次硝酸\ce{H2N2O2}及其盐}
理论上而言,\ce{N}的$+1$价含氧酸应当为\ce{H-N=O}的形式,可惜研究发现它只作为反应的不稳定中间体而存在.\ce{N}的(相对而言)稳定的$+1$价含氧酸是连二次硝酸.
\begin{substance}[\ce{H2N2O2}]
    连二次硝酸,化学式为\ce{H2N2O2},是可溶于水的不稳定白色晶体.
\end{substance}
\paragraph{\ce{H2N2O2}的结构}
与\ce{N2F2}不同的是,反式的\ce{H2N2O2}或\ce{N2O2^2-}更加稳定.
\paragraph{\ce{H2N2O2}的制备与反应}
通常通过下面的办法制备纯净的\ce{H2N2O2}:
\begin{center}
    \ce{2NaNO2 + 8Na(Hg) + 4H2O -> Na2N2O2 + 8NaOH}\\
    \ce{Na2N2O2 + 2AgNO3 -> Ag2N2O2 v + 2NaNO3}\\
    \ce{Ag2N2O2 + 2HCl ->T[\ce{Et2O}] 2AgCl v + H2N2O2}
\end{center}
最后蒸发乙醚溶剂即可.此外,还有如下办法制备\ce{Na2N2O2}:
\begin{center}
    \ce{NH2OH + RONO + 2NaOEt ->T[\ce{EtOH}] Na2N2O2 + ROH + 2EtOH}
\end{center}

\indent \ce{H2N2O2}在水溶液中是二元弱酸:
\begin{center}
    \ce{H2N2O2(aq) + H2O(l) <=> H3O+(aq) + HN2O2^-}\ \ \ $\text pK_{\text{a}1}=6.9$\\
    \ce{HN2O2^-(aq) + H2O(l) <=> H3O+(aq) + N2O2^2-}\ \ \ $\text pK_{\text{a}1}=11.6$\\
\end{center}

\indent 虽然\ce{H2N2O2}到\ce{N2}的电势高的吓人,但它在动力学上似乎不容易被还原.相反,大多数时候它可以作为还原剂,例如还原\ce{I2}:
\begin{center}
    \ce{H2N2O2 + 3I2 + 3H2O -> HNO3 + HNO2 + 6HI}
\end{center}
\ce{N2O2^2-}被氧化至\ce{N2O3^2-}有两种异构体.它们分别可以由如下方法得到:
\begin{center}
    \ce{NH2OH + BuONO2 + 2NaOMe ->T[\ce{MeOH}][$0\tc$] $\alpha$-Na2N2O3 + BuOH + 2MeOH}\\
    \ce{2Na2N2O2 + N2O4 -> 2$\beta$-Na2N2O3 + 2NO}
\end{center}
两者的结构如下所示:
\begin{figure}[H]
    \centering
    \begin{subfigure}[b]{.45\linewidth}
        \centering\includegraphics{figure/alpha-N2O32-.eps}
        \caption{\ce{$\alpha$-Na2N2O3}的结构}
    \end{subfigure}
    \begin{subfigure}[b]{.45\linewidth}
        \centering\includegraphics{figure/beta-N2O32-.eps}
        \caption{\ce{$\beta$-Na2N2O3}的结构}
    \end{subfigure}
    \caption{\ce{N2O3^2-}的两种异构体及其结构}
\end{figure}
\indent 另外,\ce{N2O2^2-}还可以作为配体.例如在\ce{[Co(NH3)5(NO)]^2+}的异构体\ce{[Co2(NH3)10(N2O2)]^4+}中就存在\ce{N2O2^2-}作为桥连配体,其结构如下:
\chemfig{Co2N2O2}{1}{\ce{[Co2(NH3)10(N2O2)]^4+}的结构}
在\ce{Pt(PPh3)2(N2O2)}中也存在\ce{N2O2^2-},其中两个\ce{O}原子对\ce{Pt}配位,形成螯合的五元环:
\chemfig{PtN2O2}{1}{\ce{Pt(PPh3)2(N2O2)}的结构}
\subsubsection{亚硝酸\ce{HNO2}及其盐}
\begin{substance}[\ce{HNO2}]
    亚硝酸,化学式为\ce{HNO2}.未分离出纯的\ce{HNO2}.
\end{substance}
\begin{substance}[\ce{NaNO2}]
    亚硝酸钠,俗称硝精,化学式为\ce{NaNO2},为淡黄色至无色的晶体,熔点为$284\tc$,沸点为$320\tc$.
\end{substance}
\paragraph{\ce{HNO2}及其盐的制备}
虽然没有分离出纯的\ce{HNO2},但它在气相中与\ce{H2O},\ce{NO}和\ce{NO2}存在如下平衡:
\begin{center}
    \ce{2HNO2(g) <=> H2O(g) + NO(g) + NO2(g)}\ \ \ $K=8.0\times10^5$
\end{center}
亚硝酸盐则可以通过对硝酸盐的部分还原而得到:
\begin{center}
    \ce{NaNO3 + Pb -> NaNO2 + PbO}
\end{center}
工业上采用\ce{NaOH}或\ce{Na2CO3}溶液吸收\ce{NO}和\ce{NO2}的混合物以制备\ce{NaNO2}:
\begin{center}
    \ce{Na2CO3 + NO + NO2 -> 2NaNO2 + CO2}
\end{center}
然后采用复分解反应制备其它亚硝酸盐;也可以通过复分解反应制备\ce{HNO2},例如:
\begin{center}
    \ce{Ba(NO2)2 + H2SO4 -> BaSO4 v + 2HNO2}
\end{center}
\paragraph{\ce{HNO2}及其盐的性质与应用}
\ce{HNO2}在水溶液中不稳定,可以按照下面的方式歧化分解:
\begin{center}
    \ce{3HNO2(aq) <=> H3O+(aq) + NO3-(aq) + 2NO(g)}
\end{center}
在强氧化剂的存在下,\ce{NO2-}可以被氧化为\ce{NO3-},例如可以用\ce{KMnO4}滴定\ce{NO2-}:
\begin{center}
    \ce{5NO2- + 2MnO4- + 6H+ -> 5NO3- + 2Mn^2+ + 3H2O}
\end{center}
亚硝酸盐本身也有一定氧化性.从天然盐卤提取碘单质的过程中,可以利用\ce{NaNO2}在酸性环境下与盐卤水反应:
\begin{center}
    \ce{2I- + 2NO2- + 4H+ -> I2 + 2NO + 2H2O}
\end{center}
前面已经提到过\ce{HNO2}可以与\ce{N2H5+}反应生成\ce{HN3}.后者可以与\ce{HNO2}进一步反应生成\ce{N2}与\ce{N2O}:
\begin{center}
    \ce{H\tbf{N}O2 + HN3 -> \{NNN\tbf{N}O\} + H2O -> N2 + N\tbf{N}O + H2O}
\end{center}
\indent \ce{NaNO2}有低毒性.它经常被用于处理肉类,用于防腐和增色增香.因此,要警惕食用\ce{NaNO2}含量过高的腌制肉类.\\
\indent \ce{NaNO2}在有机化学中可用于对芳香胺进行重氮化反应:
\begin{center}
    \ce{ArNH2 + NaNO2 + 2HCl -> [Ar-N#N]Cl + NaCl + 2H2O}
\end{center}
\paragraph{\ce{NO2-}的配合物}
有关于\ce{NO2-}的配合物的最重要的性质是它形成的键合异构体,即以\ce{N}配位的硝基配合物和以\ce{O}配位的亚硝酸酯配合物.一般而言,硝基配位的形式更加稳定.典型的是下面两种化合物:黄色的\ce{[Co(NH3)5(NO2)]^2+}和红色的\ce{[Co(NH3)5(ONO)]^2+}.

\subsubsection{硝酸\ce{HNO3}及其盐}
\begin{substance}[\ce{HNO3}]
    硝酸,化学式为\ce{HNO3},是无色液体,熔点$-41.6\tc$,沸点$82.6\tc$.
\end{substance}
\paragraph{\ce{HNO3}的生产与用途}
1900年以前,大规模生产\ce{HNO3}完全依赖于以下复分解反应:
\begin{center}
    \ce{2NaNO3 + H2SO4 -> Na2SO4 + 2HNO3}
\end{center}
1903年,人们发明了电弧炉中直接制取\ce{HNO3}的方法:
\begin{center}
    \ce{2N2(g) + 5O2(g) + 2H2O(l) ->T[放电] 4HNO3(l)}
\end{center}
这一方法仍然需要消耗极高的能量以提供反应的活化能\footnote{从另一个角度而言,如果该反应的活化能较低,那么大气中将不再有\ce{O2},海洋中也将充满\ce{HNO3}.},因而十分昂贵.直到人们发明了\ce{NH3}的催化氧化:
\begin{center}
    \ce{2NH3 + 5O2 ->T[\ce{Pt}] 2NO + 3H2O}\\
    \ce{2NO + O2 -> 2NO2}\\
    \ce{3NO2 + H2O -> 2HNO3 + NO}
\end{center}

\indent 硝酸在工业中有非常重要而广泛的应用.
\paragraph{无水\ce{HNO3}的性质}
无水\ce{HNO3}中发生明显的离解反应:
\begin{center}
    \ce{2HNO3 <=> H2NO3+ + NO3-}\\
    \ce{H2NO3+ <=> NO2+ + H2O}\\
    \ce{H2O + HNO3 <=> H3O+ + NO3-}
\end{center}
这导致了无水\ce{HNO3}的\ce{^15N}NMR信号中只有一个\ce{HNO3},\ce{NO2+}和\ce{NO3-}平衡混合物的混合峰.
\paragraph{硝酸盐的热分解}
硝酸盐的热分解产物主要和金属元素有关.下面是一些例子.
\begin{center}
    \ce{2NaNO3 ->T[$\Delta$] 2NaNO2 + O2}\\
    \ce{4KNO3 ->T[$\Delta$] 2K2O + 2N2 + 5O2}\\
    \ce{2Pb(NO3)2 ->T[$\Delta$] 2PbO + 4NO2 + O2}\\
    \ce{2AgNO3 ->T[$\Delta$] Ag + 2NO2 + O2}
\end{center}
稍特殊一些的是硝酸铵\ce{NH4NO3}的分解.在高温下或使用引爆剂时,\ce{NH4NO3}发生爆炸性的分解;缓慢加热时则分解生成\ce{N2O}:
\begin{center}
    \ce{2NH4NO3 ->T[$>300\tc$] 2N2 + O2 + 4H2O}\\
    \ce{NH4NO3 ->T[$200\sim260\tc$] N2O + 2H2O}
\end{center}
\paragraph{\ce{NO3-}的配合物}
大部分\ce{NO3-}的配合物中\ce{NO3-}都作为双齿配体存在,当然也有少部分例外.
\paragraph{原硝酸根\ce{NO4^3-}}
尽管不存在\ce{H3PO4}的类似物\ce{H3NO4},原硝酸根\ce{NO4^3-}的盐却是已经制得的:
\begin{center}
    \ce{NaNO3 + Na2O ->T[$300\tc$][7 days] Na3NO4}
\end{center}
其中\ce{NO4^3-}的键长异常的小.这可能是\ce{N}与\ce{O}之间显著的电荷相互作用所致.
\end{document}